% Page-1 teaser. Contract: references/paper-writing-template.md section 6.
%
% Tier C structure (TikZ, no external asset) carrying Tier B numbers from
% paper/data/teaser.tex. Carries NO \caption and no float number: its own
% labels do the work. Numbers are never typed here.
%
% Reading path: question -> design that answers it -> supported takeaway.
% The legend strip at the bottom is the visual key every later figure reuses.
%
% Layout is content-length-robust by construction, not by coordinates tuned
% for one paper's numbers:
%   - the Finding box and the factor matrix size to their own content
%     (TikZ `fit`/`positioning`), and every element below them measures that
%     actual extent (via `let`) instead of assuming a fixed height;
%   - the factor matrix is built from \heroArms below and accepts any number
%     of rows -- add or remove a row there and the border, the caption, and
%     the legend band all move with it automatically;
%   - row labels get a bounded `text width` (a cap, not the old unconstrained
%     single line) so a long method name wraps instead of running into the
%     cell grid.
% Extending \heroArms or writing a longer \TSfinding therefore grows this
% figure; it does not overlap or clip. Keep this discipline when editing.
%
% source-data: paper/data/teaser.tex
% generator:   scripts/make_paper_data.py --teaser
\begin{center}
  \resizebox{0.99\linewidth}{!}{%
  \begin{tikzpicture}[x=1cm,y=1cm,font=\sffamily]
    % ---- band 1: the question ------------------------------------------
    \node[anchor=west,font=\bfseries\small,text=ACBlue] at (0.42,5.12) {Question};
    \node[anchor=west,font=\small,text=ACInk,text width=10.8cm,align=left]
      at (1.95,5.12) {\TSquestion};
    \draw[ACMuted!35,line width=0.5pt] (0.42,4.72) -- (12.98,4.72);

    % ---- column titles ---------------------------------------------------
    \node[anchor=base west,font=\bfseries\small,text=ACBlue] at (0.42,4.46) {Design};
    \node[anchor=base,font=\bfseries\small,text=ACBlue] at (6.55,4.46)
      {Factors held apart};
    \node[anchor=base west,font=\bfseries\small,text=ACBlue] at (10.05,4.46) {Finding};

    % ---- left column: the registered design ------------------------------
    \node[acbox,anchor=north west,text width=2.45cm,minimum height=0.86cm,align=left,
          font=\scriptsize] (herodesign1) at (0.42,4.20)
      {\textbf{\TSsplitseeds} split seeds\\registered in advance};
    \node[acbox,anchor=north west,text width=2.45cm,minimum height=0.86cm,align=left,
          font=\scriptsize] (herodesign2) at (0.42,3.28)
      {\textbf{\TStrajper} trajectories\\per split and method};
    \node[acsoft,anchor=north west,text width=2.45cm,minimum height=0.86cm,align=left,
          font=\scriptsize] (herodesign3) at (0.42,2.36)
      {\textbf{\TSruns} preserved runs\\one manifest each};

    % ---- centre: the factor matrix ---------------------------------------
    % Structure only. Cells carry factor identity, not values, so nothing here
    % can be read as a magnitude that was never measured.
    %
    % \heroArms is a label/tint pair per row. Add a row here for a fourth or
    % fifth comparison arm; nothing else in this file needs to change. Each
    % row's vertical slot is measured from the label actually typeset in it
    % (`let` against the node's own north/south), not assumed to be one line:
    % a name too long for 2.35cm wraps to two lines and the row grows to fit
    % it, so a wrapped label can never clip into the row above, the row
    % below, or the matrix border.
    \def\heroArms{{\methodcontrol}/ACCompare,{\methodmax}/ACPrimary,{\methodrandom}/ACPrimary}
    \pgfmathsetlengthmacro{\heroMatrixTop}{4.20cm}
    \pgfmathsetlengthmacro{\heroCellH}{0.40cm}
    \pgfmathsetlengthmacro{\heroRowGap}{0.18cm}
    \pgfmathsetlengthmacro{\heroCursorY}{\heroMatrixTop-0.52cm}
    \foreach \heroLab/\heroTint [count=\heroI] in \heroArms {
      % text width is a cap: a name longer than fits on one line wraps here,
      % never past the label column and into the cell grid.
      \node[anchor=north west,font=\scriptsize,text=ACInk,text width=2.35cm,align=left]
        (herorowlab-\heroI) at (3.72,\heroCursorY) {\heroLab};
      \path let \p1=(herorowlab-\heroI.north), \p2=(herorowlab-\heroI.south)
        in \pgfextra{\xdef\heroLabTopRaw{\y1}\xdef\heroLabBotRaw{\y2}};
      \pgfmathsetlengthmacro{\heroLabH}{\heroLabTopRaw-\heroLabBotRaw}
      \pgfmathsetlengthmacro{\heroSlotH}{max(\heroLabH,\heroCellH)}
      \pgfmathsetlengthmacro{\heroCellTop}{\heroCursorY-(\heroSlotH-\heroCellH)/2}
      \pgfmathsetlengthmacro{\heroCellBot}{\heroCellTop-\heroCellH}
      \foreach \heroC in {0,1,2,3,4}{
        \fill[\heroTint,rounded corners=1.2pt]
          (6.30+\heroC*0.60,\heroCellBot) rectangle +(0.48,\heroCellH);
        \draw[ACMuted!45,line width=0.35pt,rounded corners=1.2pt]
          (6.30+\heroC*0.60,\heroCellBot) rectangle +(0.48,\heroCellH);
      }
      % \pgfmathsetlengthmacro defines \heroCursorY inside this iteration's
      % own group; re-exporting it with \global\let is what lets the next
      % iteration see the advanced value instead of the loop's start value.
      \pgfmathsetlengthmacro{\heroCursorY}{\heroCursorY-\heroSlotH-\heroRowGap}
      \global\let\heroCursorY\heroCursorY
    }
    \pgfmathsetlengthmacro{\heroLastRowBot}{\heroCursorY+\heroRowGap}
    \coordinate (heromatrixtl) at (3.55,\heroMatrixTop);
    \coordinate (heromatrixbr) at (9.60,\heroLastRowBot);
    \node[fit=(heromatrixtl)(heromatrixbr),draw=ACBlue!45,line width=0.8pt,
          rounded corners=3pt,inner sep=0pt] (heromatrixbox) {};
    % Every offset below is written with an explicit `cm` because it is
    % combined, via subtraction, with a `pgfmathsetlengthmacro` result: pgf
    % arithmetic between an explicit-unit length and a bare number treats
    % the bare number as bare points, not as a multiple of the picture's
    % x=1cm/y=1cm basis -- a dropped `cm` here silently shrinks a 0.36cm gap
    % to 0.36pt. (This is exactly the bug an earlier draft of this file hit:
    % the split-seed axis label collapsed onto the header above it.)
    \node[anchor=south,font=\scriptsize,text=ACMuted] at (7.60,\heroMatrixTop-0.36cm) {\splitseed};
    \node[anchor=north,font=\scriptsize,text=ACMuted,text width=6.4cm,align=center]
      (heromatrixcaption) at (6.55,\heroLastRowBot-0.16cm)
      {each cell aggregates every \trajseed\ under one split};
    \draw[acarrow] (3.14,3.15) -- (3.51,3.15);

    % ---- right column: the supported takeaway ----------------------------
    % No minimum height on the Finding box: it sizes to \TSfinding, however
    % long the supported sentence is. The delta box is anchored below it
    % (`positioning`'s below=of), so the two can never overlap regardless of
    % how tall the finding grows.
    \node[acsoft,anchor=north west,text width=2.68cm,minimum height=0.86cm,align=left,
          font=\scriptsize] (herofinding) at (10.05,4.20) {\TSfinding};
    \node[acbox,text width=2.68cm,minimum height=0.85cm,align=left,font=\scriptsize,
          below=0.15cm of herofinding] (herodelta)
      {\textbf{\TSdelta} $\pm$ \TSdeltaunc\\\TSmetric, lower is better};
    \draw[acarrow] (9.64,3.15) -- (10.01,3.15);

    % ---- everything below the three columns is measured, not assumed -----
    % The lowest of {left design column, matrix + its caption, finding +
    % delta} sets where the legend band starts. Whichever column is tallest
    % this time is the one that governs; the other two just have slack.
    \node[fit=(herodesign3)(heromatrixcaption)(herodelta),inner sep=0pt]
      (herocontentbottom) {};
    \path let \p1 = (herocontentbottom.south)
      in \pgfextra{\xdef\heroContentBottomRaw{\y1}};
    \pgfmathsetlengthmacro{\heroLegendLineY}{\heroContentBottomRaw-0.20cm}
    \pgfmathsetlengthmacro{\heroSwatchTopY}{\heroLegendLineY-0.32cm}
    \pgfmathsetlengthmacro{\heroSwatchBotY}{\heroSwatchTopY-0.30cm}
    \pgfmathsetlengthmacro{\heroFrameBotY}{\heroSwatchBotY-0.68cm}
    \pgfmathsetlengthmacro{\heroFrameInnerBotY}{\heroFrameBotY+0.14cm}

    % ---- band 3: the visual key later figures reuse ----------------------
    \draw[ACMuted!35,line width=0.5pt] (0.42,\heroLegendLineY) -- (12.98,\heroLegendLineY);
    \foreach \x/\tint/\lab in {%
      0.42/ACCompare/{\methodcontrol},
      4.60/ACPrimary/{compared selection methods},
      9.30/ACNeutral/{secondary configurations}}{
      \fill[\tint,rounded corners=1.2pt] (\x,\heroSwatchBotY) rectangle +(0.42,0.30);
      \draw[ACMuted!50,line width=0.35pt,rounded corners=1.2pt]
        (\x,\heroSwatchBotY) rectangle +(0.42,0.30);
      \node[anchor=west,font=\scriptsize,text=ACMuted] at (\x+0.56,\heroSwatchTopY-0.15cm) {\lab};
    }

    % ---- frame ------------------------------------------------------------
    % Drawn last but placed behind everything (background layer), sized to
    % whatever the content above actually needed. In the common case (three
    % rows, a one-line finding) this reproduces the original fixed 0..5.6
    % card exactly; it only grows when a row or a finding needs the room.
    \begin{scope}[on background layer]
      \fill[ACPrimary!35,rounded corners=4pt] (0,\heroFrameBotY) rectangle (13.4,5.6);
      \fill[white,rounded corners=4pt] (0.14,\heroFrameInnerBotY) rectangle (13.26,5.46);
    \end{scope}
  \end{tikzpicture}%
  }
\end{center}
